%% ma: L11-B16-0097
%% lop: 11
%% chuong: 5
%% bai: 16
%% dang: 11.16.1
%% loai: TL
%% muc: VDC
%% dap_an: "-"
%% nguon: "Ảnh giáo viên gửi 25/11/2026: Giới hạn hàm số dạng 0/0 (3 bài, có ĐS), Bài 3 câu 5"
%% kien_thuc: "Lớp 11 Bài 16 (giới hạn của hàm số); hằng đẳng thức x^k-1=(x-1)(x^(k-1)+...+1); tổng 1+2+...+n (Lớp 11 Bài 6 cấp số cộng)"
%% trang_thai: chinh-thuc
%% ghi_chu: "Đề gốc ghi m, n là số nguyên, đã sửa thành số nguyên dương."
\begin{ex}
Tính giới hạn $A=\lim\limits_{x\to1}\dfrac{x+x^2+x^3+\dots+x^n-n}{x+x^2+x^3+\dots+x^m-m}\ (m,n\text{ là số nguyên dương})$.
\loigiai{Tử $=(x-1)+(x^2-1)+\dots+(x^n-1)=(x-1)\left[1+(x+1)+\dots+(x^{n-1}+\dots+1)\right]$, tương tự với mẫu. Vậy $A=\dfrac{1+2+\dots+n}{1+2+\dots+m}=\dfrac{n(n+1)}{m(m+1)}$.}
\end{ex}
